% 阅读副本；权威源码为 01_source/chapters/ch01a.tex—ch01c.tex。
% 原书第 1—5 页；章首导言、第 1.1—1.2 节。
\chapter{统计模型与因果模型}
\label{ch:statistical-causal-models}

借助\textbf{统计学习（statistical learning）}，我们试图从\textbf{观测数据（observational data）}中推断随机变量之间相依关系的性质。例如，根据两个随机变量的一组联合观测样本，我们可以构造一个预测器：当获得其中一个变量的新取值时，用它来较好地估计另一个变量。这类预测所依据的理论已经相当成熟。尽管这些理论通常针对较为简单的情形，却已为我们理解如何从数据中学习提供了深刻认识。

本章将介绍其中一些认识，主要有两个原因。第一，它们有助于我们体会\textbf{因果推断（causal inference）}究竟难在哪里。在因果推断中，基础模型不再是随机变量的某个固定联合分布，而是一种能够导出多个联合分布的结构。第二，尽管关于因果估计的有限样本结果还比较少，我们仍应记住：进入更复杂的因果问题后，基本的统计估计问题并不会消失。与纯粹的统计学习中不会出现的那些因果难题相比，这些统计问题也许显得次要，但仍然需要处理。在这些准备工作的基础上，本章还将借助两个例子，循序渐进地介绍因果性的基本概念；其中一个例子是机器学习领域所熟悉的。

\section{概率论与统计学}
\label{sec:probability-statistics}

概率论与统计学以\textbf{随机试验（random experiment）}或\textbf{概率空间（probability space）} $(\Omega,\mathcal{F},P)$ 为基本模型。其中，$\Omega$ 是包含所有可能结果的集合，$\mathcal{F}$ 是由事件 $A\subseteq\Omega$ 构成的集合族，而 $P$ 是为每个事件赋予概率的测度。给定这一数学结构，概率论使我们能够对随机试验的结果进行推理。统计学习所处理的则基本上是相反的问题：已知试验结果，希望据此推断其背后数学结构的性质。例如，假设我们已经观测到数据
\begin{equation}
 (x_1,y_1),\ldots,(x_n,y_n),
 \tag{1.1}
\end{equation}
其中，$x_i\in\mathcal{X}$ 是\textbf{输入（input）}，有时也称为\textbf{协变量（covariate）}或样例（case）；$y_i\in\mathcal{Y}$ 是\textbf{输出（output）}，有时也称为目标值（target）或\textbf{标签（label）}。我们可以假设，各个 $(x_i,y_i)$（$i=1,\ldots,n$）都由同一个未知随机试验独立生成。更准确地说，这一模型假定，观测值 $(x_1,y_1),\ldots,(x_n,y_n)$ 是随机变量 $(X_1,Y_1),\ldots,(X_n,Y_n)$ 的实现值；这些随机变量\textbf{独立同分布（independent and identically distributed，i.i.d.）}，共同服从\textbf{联合分布（joint distribution）} $P_{X,Y}$。这里，$X$ 和 $Y$ 分别是取值于度量空间 $\mathcal{X}$ 和 $\mathcal{Y}$ 的随机变量。\footnote{随机变量 $X$ 是一个可测函数 $\Omega\to\mathcal{X}$，其中度量空间 $\mathcal{X}$ 配备了 Borel $\sigma$-代数。$X$ 在 $\mathcal{X}$ 上的分布 $P_X$，可由底层概率空间 $(\Omega,\mathcal{F},P)$ 中的测度 $P$ 导出。我们不必始终考虑这个底层空间，通常可以直接从随机变量的分布出发，假定随机试验直接给出从该分布抽取的数值。}
统计学和机器学习几乎都以独立同分布的数据为基础。然而，实际中的独立同分布假设可能以多种方式被破坏，例如发生\textbf{分布偏移（distribution shift）}，或系统受到\textbf{干预（intervention）}。我们将在后文看到，其中有些情形与因果性有着密切联系。

我们可能关心 $P_{X,Y}$ 的某些性质，例如：
\begin{enumerate}[label=(\roman*)]
 \item 给定输入时输出的期望，即 $f(x)=\mathbb{E}[Y\mid X=x]$。这称为\textbf{回归（regression）}，其中通常有 $\mathcal{Y}=\mathbb{R}$。
 \item 将每个 $x$ 归入概率较大类别的\textbf{二分类器（binary classifier）}，即
 \[
   f(x)=\argmax_{y\in\mathcal{Y}}P(Y=y\mid X=x),
   \qquad \mathcal{Y}=\{\pm1\}.
 \]
 \item $P_{X,Y}$ 的密度 $p_{X,Y}$，假定这一密度存在。
\end{enumerate}
在实际应用中，我们希望从有限数据集中估计这些性质。也就是说，我们依据样本 (1.1) 进行估计；等价地，也可以依据\textbf{经验分布（empirical distribution）} $P^n_{X,Y}$ 进行估计，这一分布在每个观测值上放置等权重的点质量。

这构成了一个\textbf{逆问题（inverse problem）}：我们希望估计某个无法直接观测的对象——这里是数据背后的分布——的某种性质，而所依据的观测值，是对这个对象施加某种操作后得到的。在这里，这一操作就是从未知分布中抽样。

\section{学习理论}
\label{sec:learning-theory}

现在，仿照从 $P_{X,Y}$ 得到 $f$ 的做法，假设我们用经验分布来构造经验估计 $f^n$。这实际上是一个\textbf{不适定问题（ill-posed problem）}［例如 Vapnik, 1998］，因为对样本 $(x_1,y_1),\ldots,(x_n,y_n)$ 中没有出现过的任何 $x$，相应的条件期望都没有定义。当然，我们可以先在观测样本上定义函数 $f$，再按照某种固定规则将其延拓到样本之外，例如在样本之外令 $f$ 等于 $+1$，或者选择一个连续的分段线性函数。然而，无论选择哪一种规则，输入——即经验分布——的微小变化，都可能引起输出的巨大变化。无论观测数量有多大，经验分布通常都不能完美地逼近真实分布；这种逼近中的微小误差，因而可能在估计中造成很大的误差。这意味着，如果不对经验估计 $f^n$ 所属的函数类作出额外假设，我们就无法保证这些估计能够在适当的意义下逼近最优目标 $f$。

统计学习理论借助\textbf{容量度量（capacity measure）}来形式化这些假设。如果所用函数类非常丰富，足以拟合几乎所有可以设想的数据集，那么能够拟合手头的数据并不令人意外。反之，如果事先将函数类的容量限制得很小，那么在所有可能的数据集中，只有少数数据集能够由该函数类中的某个函数来解释。在这种情况下，如果我们仍然能够解释手头的数据，就有理由相信，我们找到了蕴藏在数据中的某种规律。此时，对于未来从同一分布 $P_{X,Y}$ 抽取的数据，我们可以为解的准确性提供概率保证。

也可以从另一个角度理解这一点：我们所用的函数类融入了与观测数据背后规律相吻合的\textbf{先验知识（a priori knowledge）}，例如函数的光滑性。这类知识可以通过多种方式引入，而不同机器学习方法的差异之一，就在于如何处理这个问题。在贝叶斯方法中，我们为函数类和噪声模型指定先验分布；在正则化理论中，我们构造适当的正则化项，并将其纳入优化问题，以引导解的选择。

统计学习的困难，主要在于我们试图依据经验数据求解一个逆问题。如果已知完整的概率模型，这些问题便不复存在。讨论因果模型时，我们将看到，因果学习在某种意义上更加困难，因为它存在两个层面的不适定性。第一层是统计意义上的不适定性：无论有限样本有多大，它都不可能包含底层分布的全部信息。第二层不适定性则来自这样一个事实：即使完全知道观测分布，通常仍不能确定其背后的因果模型。

下面更详细地考察统计学习问题，并聚焦于\textbf{二类模式识别（binary pattern recognition）}，也就是二分类问题［例如 Vapnik, 1998］。此时，$\mathcal{Y}=\{\pm1\}$。我们希望根据从未知分布 $P_{X,Y}$ 独立同分布抽取的观测样本 (1.1)，学习函数 $f\colon\mathcal{X}\to\mathcal{Y}$。我们的目标是在某个函数类 $\mathcal{F}$ 中，使期望误差或\textbf{风险（risk）}\footnote{这里的“风险”概念来自统计学习理论［Vapnik, 1998］，其含义并不总与日常用法相同；它的源头可以追溯到统计决策理论［Wald, 1950；Ferguson, 1967；Berger, 1985］。在这一语境下，$f(x)$ 被理解为观测到 $x$ 后采取的行动，而损失函数衡量当真实状态为 $y$ 时，这一行动所造成的损失。}
\begin{equation}
 R[f]=\int \frac{1}{2}\lvert f(x)-y\rvert\,\mathrm{d}P_{X,Y}(x,y)
 \tag{1.2}
\end{equation}
最小。注意，这里是关于测度 $P_{X,Y}$ 的积分；如果 $P_{X,Y}$ 相对于 Lebesgue 测度存在密度 $p(x,y)$，那么这一积分就可以写成
\[
 \int \frac{1}{2}\lvert f(x)-y\rvert\,p(x,y)\,\mathrm{d}x\,\mathrm{d}y.
\]
由于 $P_{X,Y}$ 未知，我们无法计算式 (1.2)，更不用说将其最小化。因此，我们转而采用某种\textbf{归纳原则（induction principle）}，例如\textbf{经验风险最小化（empirical risk minimization）}。具体做法是在 $f\in\mathcal{F}$ 中，选择使训练误差或\textbf{经验风险（empirical risk）}
\begin{equation}
 R^n_{\mathrm{emp}}[f]
 =\frac{1}{n}\sum_{i=1}^{n}\frac{1}{2}\lvert f(x_i)-y_i\rvert
 \tag{1.3}
\end{equation}
最小的函数。

从渐近角度看，一个重要问题是，这样的学习过程是否具有\textbf{一致性（consistency）}。其基本含义是：随着 $n$ 趋于无穷，所产生的函数序列的风险，依概率收敛于给定函数类 $\mathcal{F}$ 中能够达到的最小风险。附录 A.3 将说明，只有当函数类足够“小”时，这一点才可能成立。\textbf{Vapnik--Chervonenkis（VC）维}［Vapnik, 1998］是度量函数类容量或大小的一种方式。它还能帮助我们得到有限样本保证：以较高概率，风险 (1.2) 不超过经验风险再加上一个随函数类 $\mathcal{F}$ 的大小而增大的项。

这一理论与已有的\textbf{普适一致性（universal consistency）}结果并不矛盾。普适一致性是指，学习算法的风险收敛于所有函数中能够达到的最低风险。确实存在具有普适一致性的学习算法，例如最近邻分类器和支持向量机［Devroye et al., 1996；Vapnik, 1998；Schölkopf and Smola, 2002；Steinwart and Christmann, 2008］。普适一致性在本质上告诉我们，在数据量趋于无穷的极限下，任何问题都可以学习；但它并不意味着，每个问题都能凭借有限数据得到很好的学习结果，因为还可能出现\textbf{慢收敛现象（slow rates）}。对于任意学习算法，总存在某些问题，使其学习收敛速率任意缓慢［Devroye et al., 1996］。不过，普适一致性确实表明，只要固定分布并收集足够多的数据，最终就能够任意逼近最低风险。

在实际应用中，机器学习系统近年来的成功似乎表明，有时我们确实已经进入上述渐近情形，并且往往取得了出色结果。人们投入了大量思考，力求设计出尽可能高效利用数据的方法，以便从给定数据集中得到尽可能好的结果；同时，也投入了大量精力构建可用于训练这些方法的大规模数据集。然而，在所有这些情形下，一个关键前提是：底层分布在训练与测试之间不能因干预或其他变化而发生改变。本书将论证，如果只是用概率分布描述底层规律，而不引入额外结构，我们就缺乏恰当的手段来描述究竟哪些部分可能发生变化。


% 原书第 5—7 页；第 1.3 节、图 1.1—1.2、原理 1.1。
\FloatBarrier
\section{因果建模与学习}
\label{sec:causal-modeling-learning}

\textbf{因果建模（causal modeling）}从另一种可以说更为根本的结构出发。一个因果结构能够导出一个概率模型，但它还包含后者所没有的额外信息，参见第 1.4 节的例子。按照本书的术语约定，\textbf{因果推理（causal reasoning）}是指从因果模型中推出结论的过程，这类似于概率论使我们能够对随机试验的结果进行推理。然而，由于因果模型包含的信息多于概率模型，因果推理也就比概率推理更有力：它使我们能够分析干预或分布变化所产生的影响。

正如统计学习是概率论的逆问题，我们也可以考虑，如何根据因果结构在经验层面上的表现来推断因果结构本身。这里的经验信息可以完全来自观测，也可以包括干预下的数据，例如随机化试验数据，或者分布变化下的数据。研究者用不同术语来指称这类问题，包括\textbf{结构学习（structure learning）}和\textbf{因果发现（causal discovery）}。另一个密切相关的问题是：原则上，因果结构的哪些部分能够从联合分布中推断出来？我们将其称为\textbf{结构可识别性（structure identifiability）}问题。与第 1.2 节介绍的标准统计学习问题不同，即使完全知道分布 $P$，这个问题也不会变得容易，我们仍需作出额外假设，参见第 2、4、7 章。

\input{figures/figure_1_1}

不过，这一困难不应使我们忽略通常统计问题中的不适定性依然存在。因此，在因果问题中，同样有必要关注函数类的容量，例如使用加性噪声模型，见第 4.1.4 节。只不过，我们现在试图估计的是比概率结构更加丰富的结构，因而又多了一层困难。我们将这一整体问题称为\textbf{因果学习（causal learning）}。图 1.1 总结了上述问题和模型之间的关系。

要从观测分布中学习因果结构，我们需要理解因果模型与统计模型之间的联系。第 4 章和第 7 章将重新讨论这一问题，这里先给出一个例子。人们熟知的一个说法是：“相关并不意味着因果关系。”换言之，仅靠统计性质不能确定因果结构。另一个较少为人所知的观点则是：虽然我们不能推断出具体的因果结构，但至少可以从统计相依关系推断出因果联系的存在。Reichenbach~[1956] 最先认识到了这一点。下面将他的这一认识表述为一个原理，另见图 1.2。\footnote{为使表述清晰，我们将某些重要假设称为原理。本书并不始终将这些假设视为当然成立；在这个意义上，它们不是公理。}

\input{figures/figure_1_2}

\begin{principle}[Reichenbach 共同原因原理]
\label{principle:1.1}
如果两个随机变量 $X$ 与 $Y$ 统计相依（$X\nindep Y$），那么就存在第三个变量 $Z$，对两者都产生因果影响。作为特殊情形，$Z$ 可以与 $X$ 或 $Y$ 中的一个重合。此外，变量 $Z$ 能够将 $X$ 与 $Y$ 彼此\textbf{屏蔽（screening off）}：给定 $Z$ 后，两者独立，即 $X\indep Y\mid Z$。
\end{principle}

在实际问题中，相依关系也可能由于共同原因原理未提及的其他原因而出现，例如：
\begin{enumerate}[label=(\arabic*)]
 \item 我们观测到的随机变量，实际上已经以其他变量为条件；这种条件化往往由选择偏倚隐含地引入。我们将在后文重新讨论这一问题，参见注 6.29。
 \item 随机变量只是\emph{看起来}相依。例如，我们可能在大量随机变量对中进行搜索，却没有进行多重检验校正。在这种情况下，关于变量存在相依关系的推断，并不能满足预期的第一类错误控制要求，见附录 A.2。
 \item 类似地，两个随机变量也可能都带有时间依赖，并遵循某种简单的物理规律，例如指数增长。此时，这些变量看起来彼此相依；但由于独立同分布假设被破坏，应用标准的独立性检验就缺乏依据。
\end{enumerate}
尤其是在报告随机变量之间的“伪相关”时，应牢记第 (2)、(3) 点；不少大众网站都会展示这类例子。


% 原书第 7—13 页；第 1.4 节及两个小节、图 1.3—1.4、表 1.1。
\FloatBarrier
\section{两个例子}
\label{sec:two-examples}

\subsection{模式识别}
\label{sec:pattern-recognition}

第一个例子是\textbf{光学字符识别（optical character recognition）}，这是机器学习中已经得到充分研究的问题。它并不是介绍因果结构时通常会采用的例子，却可能有助于熟悉机器学习的读者理解相关概念。下面我们将描述两个因果模型，它们都能在两个随机变量之间产生相依关系。在这个例子中，我们把这两个变量分别记为手写数字图像 $X$ 和\textbf{类别标签（class label）} $Y$。两个模型的底层因果结构不同，却会导出相同的统计结构。

\textbf{模型（i）}假设，我们通过如下过程生成每一对观测值：向书写者依次提供类别标签 $y$，并要求其每次都画出与标签对应的手写数字图像 $x$。我们假定书写者会尽力正确完成任务，但其感知类别标签的过程，以及执行绘图动作的过程，都可能存在噪声。我们可以把图像 $X$ 写成类别标签 $Y$ 和某个独立噪声 $N_X$ 的函数 $f$，并将这一函数理解为生成图像的\textbf{机制（mechanism）}。这里，$Y$ 也被视为随机变量，见图 1.3 左图。于是，我们可以根据 $P_Y$、$P_{N_X}$ 和 $f$ 计算 $P_{X,Y}$。这一分布称为\textbf{观测分布（observational distribution）}；“观测”意味着，我们只是被动地观察系统，而没有进行干预。$X$ 和 $Y$ 是相依的随机变量，因此，我们能够从观测中学习由 $x$ 到 $y$ 的映射，并且根据图像 $x$ 预测正确标签 $y$，其准确性优于随机猜测。

在这一因果结构中，有两类可能的干预，分别会导出相应的\textbf{干预分布（intervention distribution）}。\footnote{第 6.3 节将说明，更一般地说，干预可以理解为对函数和随机变量的改变。}
如果干预生成后的图像 $X$，例如修改图像，或者在图像生成后将其替换为另一幅图像，那么，先前提供给书写者并记录在数据集中的类别标签不会受到影响。形式上，改变 $X$ 不会影响 $Y$，因为 $Y\coloneqq N_Y$。相反，对 $Y$ 进行干预，意味着改变提供给书写者的类别标签，这会显著影响生成的图像。形式上，改变 $Y$ 会影响 $X$，因为 $X\coloneqq f(Y,N_X)$。图中的箭头体现了这种方向性，我们用它表示\textbf{直接因果关系（direct causation）}。

在另一种\textbf{模型（ii）}中，我们不向书写者提供类别标签，而是让书写者自行决定要写哪个数字，并同时记录相应的类别标签。在这种情况下，图像 $X$ 和记录下来的类别标签 $Y$ 都是书写者意图的函数。我们将这一书写意图记为 $Z$，同样把它看作随机变量。为使模型具有一般性，我们不仅假设图像生成过程含有噪声，也假设类别标签的记录过程含有噪声；这些噪声项仍然相互独立，见图 1.3 右图。需要注意，只要适当选择函数和噪声项，就能使这一模型导出的观测分布 $P_{X,Y}$ 与模型（i）的观测分布完全相同。\footnote{事实上，命题 4.1 表明，任意联合分布 $P_{X,Y}$ 都可以由这两个模型分别导出。}

\input{figures/figure_1_3}

现在考察模型（ii）中可能的干预。如果干预图像 $X$，情况与刚才相同，类别标签 $Y$ 不受影响。但是，如果干预类别标签 $Y$，即修改书写者记录下来的标签，那么与前一种模型不同，这次并不会影响图像。

总之，如果不对所涉及的函数类和分布类施加限制，那么模型（i）与模型（ii）可以在 $X$ 和 $Y$ 上诱导出相同的观测分布，却产生不同的干预分布。在一切都由 $P_{X,Y}$ 导出的纯概率描述中，我们看不到这种差别。然而，只要纳入关于 $P_{X,Y}$ 如何产生的结构性知识，特别是图结构、函数和噪声项，就能够讨论这一差别。

模型（i）和模型（ii）都是\textbf{结构因果模型（structural causal model，SCM）}的例子；这类模型有时也称为\textbf{结构方程模型（structural equation model）}［例如 Aldrich, 1989；Hoover, 2008；Pearl, 2009；Pearl et al., 2016］。在结构因果模型中，所有相依关系都由函数生成：这些函数根据其他变量计算某个变量的值。关键在于，应把这些函数读作\textbf{赋值（assignment）}，也就是按计算机科学中的函数来理解，而不能仅将它们看作数学等式。我们通常把这些函数视为对物理机制的建模。一个结构因果模型会导出所有可观测变量的联合分布。我们已经看到，同一个分布可以由不同的结构因果模型生成，因此，从结构因果模型转向相应的概率模型时，有关干预效应的信息，以及第 6.4 节将讨论的\textbf{反事实（counterfactual）}信息，都可能丢失。本书以结构因果模型为出发点，并尝试在这一基础上展开后续内容。

最后，就这个例子再说明两点。

第一，图 1.3 很好地展示了 Reichenbach 的共同原因原理。$X$ 和 $Y$ 的相依关系可以有多种因果解释。在右图中，给定 $Z$ 后，$X$ 与 $Y$ 就相互独立：图像与标签之间不存在书写意图之外的共享信息。

第二，有时人们认为，只有把时间概念纳入考虑，才能讨论因果性。时间确实在上述例子中发挥了作用，例如，它使我们能够排除“干预 $X$ 会影响类别标签”这种可能性。不过，这并没有什么问题：统计数据集由某个随时间展开的过程生成，本来就是很常见的情形。例如，在模型（i）中，$X$ 与 $Y$ 之间的统计相依关系，源于一个动态过程。书写者先读取标签、规划动作；这涉及大脑中的复杂活动，最终由肌肉和笔执行动作。我们对这一过程的理解还不完整，但它确实是一个随时间展开的物理动态过程，其最终结果是 $X$ 和 $Y$ 之间一个非平凡的联合分布。进行统计学习时，我们只关心这一最终结果。

因此，不仅因果结构，纯粹的概率结构也可以由随时间展开的过程产生；甚至可以认为，这归根结底是它们唯一可能的产生方式。不过，在这两类描述中，暂时不考虑时间往往都有助于分析。在统计学中，讨论统计相依等概念时，通常不必涉及时间；在因果模型中，讨论干预效应时，也往往不必涉及时间。这两个描述层次都可以看作对底层物理模型的抽象；后者更加精确，对现实的描述也比它们更完整，见表 1.1。此外，模型中的变量未必对应明确界定的时点。例如，心理学家研究学生学习动机与学习表现之间的统计关系或因果关系时，就很难为这两个变量分别指定某个具体时点。对应明确时点的测量，反而更常见于物理学、化学等“硬科学”。

\input{figures/table_1_1}

\FloatBarrier
\subsection{基因扰动}
\label{sec:gene-perturbation}

第 1.4.1 节已经说明，不同的因果结构会产生不同的干预分布。有时，我们确实希望预测某个随机变量在这类干预下的取值。下面考虑一个来自遗传学的例子，尽管它在某些方面作了过度简化。假设我们已经获得基因 A 的活性数据，以及与之对应的某个\textbf{表型（phenotype）}的测量值。图 1.4 左上图给出了一个示例数据集。可以看到，这两个变量高度相关。这种相关可以用于传统预测：如果观测到基因 A 的活性在 6 左右，我们就会以较高概率预期表型的取值落在 12 至 16 之间。基因 B 的情形也类似，见左下图。

另一方面，我们也可能关心敲除基因 A 后的表型，也就是将该基因的活性设为 0 后会发生什么。\footnote{为简单起见，假设我们能够获得基因的真实活性，不存在测量噪声。}
然而，在完全不了解因果结构的情况下，我们无法给出有实质内容的回答。如果基因 A 对表型具有因果影响，那么干预后，表型预计会发生显著变化，见右上图。事实上，此时甚至可能仍然使用从观测数据中学到的同一个线性模型。反之，如果存在某个共同原因，例如第三个基因 C，同时影响基因 B 的活性和表型，那么对基因 B 的干预就不会影响表型，见右下图。

与模式识别的例子类似，这里的模型也经过适当选择，使基因 A 与表型的联合分布，等于基因 B 与表型的联合分布。因此，仅凭观测数据，无论样本量是否趋于无穷，我们都无法区分图中上下两种情形。归结起来，如果不愿采用因果性方面的概念，那么，对于“敲除基因后，表型会怎样”这一预测问题，我们就只能回答：“不知道。”

\input{figures/figure_1_4}
\FloatBarrier
